%! Properties and Applications of Determinants — Unit 5, Lesson 5.2 — Group Activity
\documentclass[10pt]{article}
\usepackage{linalg-article}
\usepackage{linalg-boxes}
\newcommand{\TallMath}[1]{$\displaystyle #1\rule[-1.4em]{0pt}{3.2em}$}
\newcommand{\vv}{\mathbf{v}}
\newcommand{\ww}{\mathbf{w}}
\newcommand{\avec}{\mathbf{a}}
\newcommand{\bb}{\mathbf{b}}
\newcommand{\xx}{\mathbf{x}}
\newcommand{\T}{\mathsf{T}}
\newcommand{\zero}{\mathbf{0}}

\begin{document}

\pageheader{Unit 5, Lesson 5.2}{Group Activity}
\namepartnerperiod

\vspace{0.10in}

\begin{headlinebox}{blush}
\small\textbf{Rules, Pivots, and Products.} You no longer have to expand every determinant. Use the
\textbf{rules}: a \emph{shear} (adding a multiple of one row to another) changes \textbf{nothing}, a
\emph{swap} flips the sign, and a \emph{scale} multiplies. Reduce to triangular and $\det=\pm(\text{product
of the pivots})$. And determinants \textbf{multiply}: $\det(AB)=\det A\,\det B$. Work with your partner,
show each step, and \emph{justify} every claim.
\end{headlinebox}

\vspace{0.10in}

\begin{tcolorbox}[colback=white, colframe=black!40, title=\textbf{Tier R --- Remediate},
    fonttitle=\bfseries, arc=1.5mm, left=3mm, right=3mm, top=2mm, bottom=2mm, breakable]
Reduce to triangular, then multiply the pivots. Show your work.
\begin{enumerate}[leftmargin=1.7em, itemsep=9pt]
  \item Compute $\det\begin{bsmallmatrix} 1 & 2 & 1 \\ 2 & 5 & 3 \\ 0 & 3 & 5 \end{bsmallmatrix}$ by
        elimination (no swaps needed). List your pivots, then multiply. \writelines{2}
  \item \textbf{A shear changes nothing.} The matrix $\begin{bsmallmatrix} 3 & 1 \\ 2 & 4 \end{bsmallmatrix}$
        has determinant $10$. Add $2\times(\text{row }1)$ to row $2$, write the new matrix, and compute
        its determinant. Did it change? \; \blank{3.2cm}
  \item \textbf{A repeated column.} Compute $\det\begin{bsmallmatrix} 1 & 5 & 1 \\ 2 & 4 & 2 \\ 3 & 0 & 3 \end{bsmallmatrix}$
        (notice column~3 $=$ column~1). What does an equal pair of columns do to the box, and to $\det$? \; \blank{3.0cm}
\end{enumerate}
\end{tcolorbox}

\begin{tcolorbox}[colback=white, colframe=black!40, title=\textbf{Tier A --- Approaching Proficiency},
    fonttitle=\bfseries, arc=1.5mm, left=3mm, right=3mm, top=2mm, bottom=2mm, breakable]
Track the sign, and use the product rule.
\begin{enumerate}[leftmargin=1.7em, itemsep=9pt]
  \item \textbf{A swap flips the sign.} To start elimination on
        $A=\begin{bsmallmatrix} 0 & 1 & 2 \\ 1 & 3 & 1 \\ 2 & 1 & 0 \end{bsmallmatrix}$ you must first
        swap two rows (the top-left entry is $0$). Reduce to triangular, count your swaps, and give
        $\det A$ with the correct sign. \writelines{3}
  \item \textbf{Product rule.} Let $A=\begin{bsmallmatrix} 2 & 0 \\ 1 & 3 \end{bsmallmatrix}$ and
        $B=\begin{bsmallmatrix} 1 & 2 \\ 0 & 4 \end{bsmallmatrix}$. Compute $\det A$, $\det B$, and
        $\det A\cdot\det B$. Then multiply out $AB$ and compute $\det(AB)$ directly. Do they match? \writelines{3}
  \item \textbf{Inverses.} Using part~2, what is $\det A^{-1}$? If a matrix has $\det=0$, explain in one
        sentence why it \emph{cannot} have an inverse. \writelines{2}
\end{enumerate}
\end{tcolorbox}

\begin{tcolorbox}[colback=white, colframe=black!40, title=\textbf{Tier E --- Extension},
    fonttitle=\bfseries, arc=1.5mm, left=3mm, right=3mm, top=2mm, bottom=2mm, breakable]
\textbf{An application: Cramer's rule.} Determinants can \emph{solve} a system. For
$A\xx=\bb$, each unknown is a ratio of determinants: replace a column of $A$ with $\bb$, take the
determinant, and divide by $\det A$.
\begin{enumerate}[leftmargin=1.7em, itemsep=9pt]
  \item Solve $\begin{cases} 2x - y = 3 \\ \ x + 3y = 5 \end{cases}$ so that
        $A=\begin{bsmallmatrix} 2 & -1 \\ 1 & 3 \end{bsmallmatrix}$, $\bb=\begin{bsmallmatrix} 3 \\ 5 \end{bsmallmatrix}$.
        Compute $\det A$; then $x=\dfrac{\det A_x}{\det A}$ (replace column~1 by $\bb$) and
        $y=\dfrac{\det A_y}{\det A}$ (replace column~2 by $\bb$). \writelines{3}
  \item Check your $(x,y)$ in both original equations. \writelines{2}
  \item \textbf{Transpose rule.} Verify $\det A^{\T}=\det A$ for
        $A=\begin{bsmallmatrix} 2 & 1 \\ 4 & 3 \end{bsmallmatrix}$: compute $\det A$ and $\det A^{\T}$.
        Why does this mean every \emph{row} rule is also a \emph{column} rule? \writelines{2}
\end{enumerate}
\end{tcolorbox}

\end{document}
